\scanpage{24}
\mspage{21}

(7 juin) \unsure{Une|Un} modèle de $1$-forme au sens précédent (ensemble de lieux $L$, avec ens de ``segments'' $S\subset\mathcal P(L)$, et $\forall I\in S$, un $\delta I\in\mathcal P_2(I)$, \deleted{\unsure{de}} de façon à satisfaire les conditions (F1) à (F6)) est appelé un \uline{réseau}.

Les lieux critiques \inserted{ou singuliers} ou de branchement sont appelés aussi \uline{noeuds} du réseau. Les comp. connexes de $L_r\setminus\underset{\text{noeuds}}{\underbrace{L_s}}$ sont \unsure{appelées} les \uuline{lignes} du réseau (de préférence à arêtes). Composantes connexes : les \uline{composantes} du réseau.

On parlera de l'\uline{ordre} ou de la \uline{multiplicité} d'un noeud (plus \unsure{général$^{\mathrm t}$|généralement}, d'un lieu du réseau -- elle est égale à $2$ en un lieu \uline{régulier}).

On dira qu'un \deleted{segment} \inserted{réseau} est \uline{divisible} \unsure{ssi} pour \unsure{tt} segment $I$, $I^\circ\ne\varnothing$.

On dira qu'il est \uline{réduit} si pour \unsure{tt} segment $I$, $I^\circ=\varnothing$ \inserted{i.e. $I=\delta I$ i.e. $\operatorname{card}I=2$}, ou encore si $L_r=\varnothing$, i.e. tous ses lieux sont des noeuds.

\mspage{23}

\rotatebox{90}{\unsure{EXTRA STRONG}}
\rotatebox{90}{\unsure{EXTRA STRONG}}
